% ============================================================
% §2.2 · Fundamentos Teóricos — Manual Técnico K-QGMIP
% ============================================================

\section{Fundamentos Teóricos}

\subsection{Definición formal de k-particiones}

\subsubsection*{Notación y definición}

Sea $V = \{0, 1, \ldots, n-1\}$ el conjunto de $n$ variables binarias que componen el sistema.
Una \textbf{k-partición} de $V$ es una familia $\Pi = \{P_1, P_2, \ldots, P_k\}$ tal que:

\begin{enumerate}
  \item \textbf{No vacuidad}: $P_i \neq \emptyset$ para todo $i = 1, \ldots, k$.
  \item \textbf{Disjunción}: $P_i \cap P_j = \emptyset$ para todo $i \neq j$.
  \item \textbf{Exhaustividad}: $\bigcup_{i=1}^{k} P_i = V$.
\end{enumerate}

El caso $k = 2$ corresponde a la \textbf{bipartición} clásica $(A, B)$ donde $A \cup B = V$ y
$A \cap B = \emptyset$. El caso $k = 1$ es la \textbf{partición trivial} $\{V\}$, que no produce
ninguna pérdida de información. El caso $k = n$ es la \textbf{atomización total}, donde cada
variable forma su propia parte.

\subsubsection*{Ejemplo ilustrativo para $n = 4$}

Sea $V = \{A, B, C, D\}$ un sistema de 4 nodos. Las posibles k-particiones para distintos
valores de $k$ son:

\begin{table}[H]
  \centering
  \caption{Ejemplos de k-particiones del conjunto $V = \{A, B, C, D\}$.}
  \label{tab:ejemplos_kpart}
  \begin{tabular}{clp{7cm}}
    \toprule
    $k$ & Número de particiones $S(4,k)$ & Ejemplos representativos \\
    \midrule
    1 & 1  & $\{A,B,C,D\}$ \\
    2 & 7  & $\{A,B\} \mid \{C,D\}$;\quad $\{A\} \mid \{B,C,D\}$ \\
    3 & 6  & $\{A,B\} \mid \{C\} \mid \{D\}$;\quad $\{A\} \mid \{B,C\} \mid \{D\}$ \\
    4 & 1  & $\{A\} \mid \{B\} \mid \{C\} \mid \{D\}$ \\
    \bottomrule
  \end{tabular}
\end{table}

Para $n = 4$ hay en total $\sum_{k=1}^{4} S(4,k) = 15$ particiones distintas (el número de Bell
$B_4 = 15$). El proyecto evalúa $k \in \{2, 3, 4, 5\}$ — la partición trivial ($k=1$) se
excluye porque no produce pérdida de información.

\subsubsection*{Propiedades fundamentales}

\begin{itemize}
  \item \textbf{Refinamiento}: $\Pi'$ es un \textit{refinamiento} de $\Pi$ si cada parte de
        $\Pi'$ está contenida en alguna parte de $\Pi$. El refinamiento iterativo que implementan
        KGeoMIP y KQNodes produce familias anidadas de particiones.

  \item \textbf{Monotonicidad de la pérdida}: si $\Pi'$ refina a $\Pi$, entonces
        $\Phi^*(\Pi') \geq \Phi^*(\Pi)$. Esto se debe a que al dividir una parte en dos se
        impone una restricción adicional de independencia causal, lo que solo puede mantener
        o incrementar la pérdida.

  \item \textbf{Conteo por número de Stirling}: el número de k-particiones de un conjunto de
        $n$ elementos está dado por el \textbf{número de Stirling de segundo tipo}:
        \begin{equation}
          S(n, k) = \frac{1}{k!} \sum_{j=0}^{k} (-1)^{k-j} \binom{k}{j} j^n
          \label{eq:stirling}
        \end{equation}
        que satisface la recurrencia $S(n,k) = k \cdot S(n-1,k) + S(n-1,k-1)$.
\end{itemize}

\subsection{Formulación del problema de optimización}

\subsubsection*{Marco de IIT 4.0 — Efecto causal}

En el formalismo de la Teoría de la Información Integrada versión 4.0 (IIT 4.0), cada sistema
$\mathcal{S}$ de $n$ variables binarias se describe mediante su \textbf{Matriz de Probabilidad
de Transición (TPM)}: la matriz $\mathbf{T} \in [0,1]^{2^n \times n}$ donde $T_{s,i}$ es la
probabilidad de que la variable $i$ esté en estado ON en $t+1$, dado que el sistema estaba en
el estado $s$ en $t$.

Para un estado inicial $s_t \in \{0,1\}^n$, la \textbf{distribución marginal} del sistema
sobre el futuro es el vector:
\begin{equation}
  p(s_{t+1} \mid s_t) = \bigl(T_{s_t, 0},\; T_{s_t, 1},\; \ldots,\; T_{s_t, n-1}\bigr)
  \in [0,1]^n
  \label{eq:dist_marginal}
\end{equation}

\subsubsection*{Distribución de una k-partición}

Dada una k-partición $\Pi = \{P_1, \ldots, P_k\}$, la \textbf{distribución particionada}
asume que las partes son causalmente independientes. La distribución marginal del nodo $i$
en la parte $P_m$ se calcula marginalizando (promediando) sobre todas las variables presentes
\textbf{fuera} de $P_m$:
\begin{equation}
  p_i^{(\Pi)}(s_{t+1} \mid s_t) = \sum_{\mathbf{s}_{-P_m}} T_{s_t^{(P_m)} \otimes \mathbf{s}_{-P_m},\; i}
  \cdot p(\mathbf{s}_{-P_m})
  \label{eq:dist_particionada}
\end{equation}
donde $\mathbf{s}_{-P_m}$ representa los estados de las variables fuera de $P_m$ (que se
promedian con distribución uniforme en la implementación).

\subsubsection*{Función objetivo — EMD-Effect}

La \textbf{pérdida de información} de la k-partición $\Pi$ se mide mediante la
\textit{Earth Mover's Distance} del efecto (EMD-Effect), que en IIT 4.0 bajo la geometría
$L_1$ toma la forma analítica:
\begin{equation}
  \Phi^*(\Pi) = \mathrm{EMD}\!\left(
    p(s_{t+1} \mid s_t),\;
    \bigotimes_{m=1}^{k} p^{(P_m)}(s_{t+1} \mid s_t)
  \right)
  = \sum_{i=1}^{n} \left| p_i(s_{t+1} \mid s_t) - p_i^{(\Pi)}(s_{t+1} \mid s_t) \right|
  \label{eq:phi}
\end{equation}
donde $\bigotimes$ denota el producto tensorial (independencia causal entre partes) y la norma
$L_1$ actúa componente a componente sobre los vectores de distribución marginal.

\subsubsection*{Problema de optimización}

El problema de la \textbf{k-Partición de Mínima Información (k-MIP)} se formula como:
\begin{equation}
  \Pi^* = \arg\min_{\Pi \in \mathcal{P}(V,k)} \Phi^*(\Pi)
  \label{eq:kmip}
\end{equation}
\begin{align*}
  \text{sujeto a:} \quad
  & P_i \neq \emptyset, \quad \forall i = 1, \ldots, k \\
  & P_i \cap P_j = \emptyset, \quad \forall i \neq j \\
  & \bigcup_{i=1}^{k} P_i = V
\end{align*}
donde $\mathcal{P}(V, k)$ denota el conjunto de todas las k-particiones de $V$.

El valor óptimo $\Phi^*(\Pi^*)$ cuantifica la \textbf{información integrada de orden $k$}: si
es pequeño, el sistema puede descomponerse en $k$ partes casi independientes; si es grande,
el sistema es causalmente irreducible a nivel de $k$ partes.

\subsection{Extensión del marco teórico de bi-particiones a k-particiones}

\subsubsection*{GeoMIP: de bipartición geométrica a KGeoMIP divisivo}

GeoMIP resuelve la bipartición óptima ($k=2$) mediante un recorrido BFS sobre el
\textbf{hipercubo de Hamming} de $D$ dimensiones (donde $D$ es la dimensión del subsistema),
construyendo una tabla de costos de transición $\mathbf{T} \in \mathbb{R}^{2^D \times D}$ y
evaluando la EMD-Effect solo para los candidatos más prometedores identificados por cercanía
topológica.

Para extender a $k \geq 3$, KGeoMIP adopta la \textbf{estrategia E4 (refinamiento divisivo
top-down)}, que parte de la bipartición exacta de GeoMIP y la refina iterativamente:

\begin{enumerate}
  \item \textbf{Anclaje $k=2$}: se ejecuta GeoMIP completo para obtener la bipartición óptima
        $\Pi_2^* = \{P_1, P_2\}$. Esto garantiza $\Phi^*(\Pi_2^*) = \Phi^*_{\text{GeoMIP}}$ por
        construcción.
  \item \textbf{Construcción de la matriz de similitud causal $S$}: $S_{ij}$ cuantifica
        cuánto costo de transición acumula la variable $i$ sobre los estados donde la variable
        $j$ difiere del estado inicial (Ecuación~\ref{eq:S_matrix}).
  \item \textbf{Refinamiento iterativo}: para $t = 3, 4, \ldots, k$, se selecciona la parte
        $P^*$ de la partición actual que minimiza el \textbf{incremento de pérdida} $\Delta\Phi$
        al bipartirla, y se reemplaza $P^*$ por la bipartición de menor costo de $P^*$.
\end{enumerate}

\begin{equation}
  S_{ij} = \frac{1}{2}\sum_{s=0}^{2^D-1} T_{s,i} \cdot \mathbf{1}\!\left[
    \text{bit}_j(s) \neq \text{bit}_j(s_0)
  \right]
  + \text{(término simétrico)}
  \label{eq:S_matrix}
\end{equation}

La función $S$ guía la selección de candidatos de corte dentro de cada parte, priorizando
pares de variables con alta similitud causal (que deben quedarse juntas) frente a pares con
baja similitud (que son buenos candidatos de separación).

\subsubsection*{QNodes: de Queyranne submodular a KQNodes greedy}

QNodes resuelve la bipartición óptima ($k=2$) mediante el \textbf{algoritmo de Queyranne}
(Queyranne, 1998), que minimiza funciones submodulares simétricas en $\mathcal{O}(D^3)$
evaluaciones del oracle. La función que QNodes minimiza:
\begin{equation}
  f(A) = \sum_{i=1}^{N} \min\bigl(
    |\bar{\mu}_B(i) - v_i|,\;
    |\bar{\mu}_A(i) - v_i|
  \bigr)
  \label{eq:f_oracle}
\end{equation}
es una aproximación submodular simétrica de la EMD-Effect, donde $\bar{\mu}_A(i)$ es la media
de $T_{\cdot,i}$ sobre los estados en $A$, y $v_i = T_{s_0,i}$ es el valor pivote del nodo $i$.

Para $k \geq 3$, KQNodes aplica un \textbf{refinamiento greedy con MinHeap}: en lugar de
evaluar todas las k-particiones posibles ($S(n,k)$ candidatos), aplica Queyranne $k-1$ veces
sobre bloques progresivamente más pequeños, guiado por el criterio C4 de mínimo $\varphi$
local:

\begin{equation}
  P^* = \arg\min_{P_i \in \Pi,\; |P_i| \geq 2} \varphi_{\text{local}}(P_i)
  \label{eq:criterio_c4}
\end{equation}

\subsubsection*{Justificación de la calidad de las heurísticas}

\begin{itemize}
  \item \textbf{Exactitud para $k=2$}: ambas estrategias producen resultados exactamente
        equivalentes a sus anclas GeoMIP y QNodes para $k=2$, verificado con
        $|\Delta\Phi| < 10^{-9}$ en 13 casos de prueba.

  \item \textbf{Garantía de monotonicidad}: el refinamiento iterativo garantiza
        $\Phi^*(\Pi_{k+1}) \geq \Phi^*(\Pi_k)$ por construcción, ya que cada paso añade
        una restricción de independencia adicional.

  \item \textbf{Optimalidad dentro de cada bloque}: dentro de cada parte $P_i$, Queyranne
        garantiza la bipartición óptima bajo la función proxy $f$, que es empíricamente
        $\sim$97--100\% exacta respecto a la EMD-Effect real (Kitazono et al., \textit{Entropy}
        2018).

  \item \textbf{Gap empírico reducido}: los experimentos A/B del proyecto muestran que el
        criterio C4 de KQNodes produce un gap promedio inferior al 5\% respecto al óptimo
        exhaustivo para $n \leq 8$. KGeoMIP-E4 muestra comportamiento análogo.

  \item \textbf{Suboptimalidad global para $k \geq 3$}: ninguna de las dos estrategias garantiza
        el óptimo global. La función $f$ no es submodular para k-particiones en general, y el
        refinamiento divisivo puede quedar atrapado en mínimos locales no globales.
\end{itemize}

\subsection{Análisis de complejidad del espacio de soluciones}

\subsubsection*{Crecimiento del espacio de k-particiones}

El número de k-particiones de $n$ elementos crece superexponencialmente con $n$ y $k$.
La Tabla~\ref{tab:espacios} muestra la magnitud del espacio de búsqueda para los tamaños
de red utilizados en el proyecto:

\begin{table}[H]
  \centering
  \caption{Tamaño del espacio de soluciones $S(n,k)$ para distintas configuraciones.
           Los valores $\gg 10^6$ hacen inviable la búsqueda exhaustiva.}
  \label{tab:espacios}
  \begin{tabular}{crrrr}
    \toprule
    $n$ & $k=2$ & $k=3$ & $k=4$ & $k=5$ \\
    \midrule
    5  & 15           & 25              & 10                & 1               \\
    8  & 127          & 966             & 2\,646            & 3\,025          \\
    10 & 511          & 9\,330          & 145\,750          & 1\,082\,250     \\
    15 & 16\,383      & $\sim$2.4\,M    & \cellcolor{gray!20}$\gg 10^9$  & \cellcolor{gray!20}$\gg 10^{11}$ \\
    20 & $\sim$524\,K & \cellcolor{gray!20}$\sim$580\,M & \cellcolor{gray!20}---          & \cellcolor{gray!20}---          \\
    25 & $\sim$16\,M  & \cellcolor{gray!20}---          & \cellcolor{gray!20}---          & \cellcolor{gray!20}---          \\
    \bottomrule
  \end{tabular}
\end{table}

\subsubsection*{Comparación con la bipartición}

Para la bipartición ($k=2$), el espacio de soluciones es $S(n,2) = 2^{n-1} - 1$, que crece
\textbf{exponencialmente} con $n$. QNodes y GeoMIP reducen esto a complejidades polinomiales:
$\mathcal{O}(D^3)$ y $\mathcal{O}(D \cdot 2^D)$ respectivamente.

Para k-particiones ($k \geq 3$), el espacio crece \textbf{superexponencialmente}. La relación
asintótica de los números de Stirling para $k$ fijo es:
\begin{equation}
  S(n, k) \sim \frac{k^n}{k!} \quad \text{cuando } n \to \infty
  \label{eq:stirling_asintotico}
\end{equation}
lo que significa que el cociente $S(n,k) / S(n,2)$ crece como $k^{n-1}/k!$, haciéndose
impracticable mucho más rápido que el caso bipartición.

\subsubsection*{Reducción lograda por las heurísticas}

KGeoMIP y KQNodes transforman el problema $\mathcal{O}(S(n,k))$ en problemas polinomiales:

\begin{table}[H]
  \centering
  \caption{Reducción de complejidad respecto a la búsqueda exhaustiva.}
  \label{tab:reduccion_complejidad}
  \begin{tabular}{lll}
    \toprule
    Algoritmo & Complejidad & Factor de reducción vs. exhaustiva \\
    \midrule
    Búsqueda exhaustiva       & $\mathcal{O}(S(n,k) \cdot N)$    & $1\times$ (baseline) \\
    KQNodes (C4)              & $\mathcal{O}(k \cdot D^3 \cdot N)$ & $S(n,k) / (k \cdot D^3)$ \\
    KGeoMIP (E4)              & $\mathcal{O}(k \cdot D^2 \cdot 2^D)$ & $S(n,k) \cdot N / (k \cdot D^2 \cdot 2^D)$ \\
    \bottomrule
  \end{tabular}
\end{table}

Para $n=10$, $k=4$, $D=10$ y $N=10$: la búsqueda exhaustiva requiere $\sim$1.5 millones de
evaluaciones EMD, mientras KQNodes requiere $\leq 4 \cdot 10^3 \cdot 10 = 40\,000$ operaciones
elementales — una reducción de más de 30 veces. Para $n=20$, $k=3$, la reducción supera los
8 órdenes de magnitud.
